\documentclass[11pt]{article}
\usepackage{ap-lang}
\title{Part 1: The Rhetorical Situation}
\author{Paul Wang}
\date{November 2026}
\docpart{Part 1 \textperiodcentered\ Study Guide}

\begin{document}
\maketitle

\begin{abstract}
Every text on the exam was written by someone, for someone, at a particular moment, to get something done. This guide gives you the vocabulary for that situation: exigence and context (Section 1), audience (Section 2), purpose and the writer's persona (Section 3), rhetorical choices and the appeals (Section 4), and how writers adjust their choices to fit a situation (Section 5). Section 6 turns the same ideas into a method for the rhetorical analysis essay (Question 2). Each section ends with check questions; the key is on the last page.
\end{abstract}

\scoringline{Check questions: \blank\ / \qtotal \qquad Target: 13+ \qquad Terms I mix up: \blank[10em]}

\section{Exigence and context}

\textbf{Exigence} is the problem, event, or need that prompts the writer to write \emph{now}: a factory closing, a new law, a graduation, a public insult. \textbf{Context} is the wider set of circumstances (historical, cultural, political) in which the text appears. The exigence answers ``why this text, why now?''; the context answers ``what world does it enter?''

\begin{terms}
Writer (speaker) & The person or group producing the text, including their position and credibility: a mayor, a scientist, a former prisoner. \\
Exigence & The specific spark: ``the council's vote to close the library branch.'' \\
Context & The surrounding conditions: ``a city cutting budgets after a recession.'' \\
Message & The text's main claim or idea. \\
\end{terms}

\begin{keyidea}[A working mnemonic]
Many classes use \textbf{SPACE}: \textbf{S}peaker, \textbf{P}urpose, \textbf{A}udience, \textbf{C}ontext, \textbf{E}xigence. Fill it in from the introductory note before you read the passage.
\end{keyidea}

\begin{trap}
Exigence vs.\ context. The exigence is the trigger (a specific event); the context is the background. ``The 1906 earthquake'' is the exigence for a fundraising appeal; ``a city that had grown faster than its building codes'' is context.
\end{trap}

\begin{q}{2}{1.1 Exigence}{The specific event prompting the writing = exigence.}
A columnist writes the day after her town's only hospital announces it will close its maternity ward. The announcement is best described as the column's
\chh{audience}{exigence}{context}{purpose}
\end{q}

\begin{q}{3}{1.1 Context}{Background conditions surrounding the text = context.}
In the same column, the writer notes that rural hospitals across the state have been losing money for a decade. This information establishes the
\chh{exigence}{thesis}{context}{persona}
\end{q}

\section{Audience}

The \textbf{audience} is the group the writer wants to reach, which is not always everyone who reads the text. Writers consider the audience's \textbf{beliefs, values, needs, and background knowledge}, and especially what the audience already thinks about the subject.

\begin{compare}{Audience is\ldots}{The writer tends to\ldots}{Example}
Sympathetic & Rally and mobilize; urge action; use ``we.'' & A speech to union members before a strike. \\
Skeptical or hostile & Concede first; build common ground; lead with shared values; delay the thesis. & A letter to clergy who called the writer's protest ``unwise.'' \\
Uninformed & Define terms; explain background; use concrete examples. & A scientist explaining vaccines to parents. \\
Expert & Use technical terms; skip basics; cite data. & A memo to engineers. \\
\end{compare}

Look for audience clues: direct address (``you,'' ``my fellow graduates''), shared references, the level of explanation, and what the writer anticipates the audience will object to.

\begin{q}{1}{1.2 Audience}{Conceding and seeking common ground = writer expects resistance.}
A writer opens by praising her readers' commitment to public safety and admitting that their concerns about her proposal are ``reasonable.'' This most strongly suggests the writer expects her audience to be
\chh{skeptical}{already persuaded}{indifferent}{uninformed}
\end{q}

\begin{q}{2}{1.2 Audience}{Defining basic terms signals a non-expert audience.}
An essay on soil erosion stops to explain what ``topsoil'' is and why it matters for crops. The writer most likely assumes the audience
\ch{has read the writer's earlier work}{lacks specialized knowledge of the subject}{is made up of agricultural scientists}{already opposes conservation}
\end{q}

\section{Purpose and persona}

\textbf{Purpose} is what the writer wants the audience to think, feel, or do. It is more specific than ``to inform'' or ``to persuade'': ``to persuade the school board to delay the start of the school day.'' A text can have a primary purpose and secondary ones (to entertain on the way to persuading).

The \textbf{persona} (or \textbf{ethos}) is the character the writer presents: the patient expert, the outraged citizen, the humble newcomer, the fellow parent. Writers build credibility by showing knowledge, good judgment, fairness to opponents, and shared values with the audience.

\begin{terms}
Purpose verbs & advocate, urge, warn, defend, justify, criticize, celebrate, commemorate, reconcile, challenge, call to action, correct a misconception. \\
Ethos moves & citing experience (``in twenty years of teaching''), acknowledging limits, treating opponents fairly, invoking shared identity (``as a fellow veteran''). \\
\end{terms}

\begin{trap}
``To inform'' is rarely the answer on this exam. Passages are chosen because they argue; even a seemingly neutral history usually aims to change how the reader sees something.
\end{trap}

\begin{q}{4}{1.3 Purpose}{Purpose = what the writer wants this audience to do or believe, stated specifically.}
A commencement speaker describes her own failures and ends: ``Do not wait to be ready. No one ever is.'' Her primary purpose is to
\ch{inform graduates about her career}{entertain the audience with anecdotes}{criticize universities for poor career advice}{urge graduates to act despite uncertainty}
\end{q}

\begin{q}{4}{1.3 Ethos}{Fair treatment of opponents builds credibility.}
A writer arguing for a new bike lane states her opponents' strongest objection accurately and calls it ``a serious point.'' This choice mainly
\ch{appeals to the audience's fear}{weakens her argument by conceding defeat}{shows that she has changed her mind}{establishes her credibility as a fair-minded writer}
\end{q}

\section{Rhetorical choices and the appeals}

A \textbf{rhetorical choice} is any decision a writer makes to serve the purpose: what to include, what order to put it in, which words to use, which example to tell, whom to quote, what tone to take. The exam asks you to explain choices, so describe them in specific terms (``the writer opens with the story of a single family'') rather than labels alone.

The classical \textbf{appeals} are a useful lens:
\begin{terms}
Ethos & Appeal to the writer's credibility and character. \\
Pathos & Appeal to the audience's emotions, values, and imagination. \\
Logos & Appeal to reason: evidence, logic, cause and effect. \\
Kairos & Fitness to the moment: why this argument is timely now. \\
\end{terms}

\begin{keyidea}[Choices before labels]
``The writer uses pathos'' earns nothing on its own. ``By describing the empty swing set outside the closed daycare, the writer makes parents on the council picture their own children without care, turning a budget line into a loss they can feel'' names a choice, its effect on this audience, and its link to the purpose. That is analysis.
\end{keyidea}

\begin{q}{1}{1.4 Choices}{Analysis names the specific choice and its effect on the audience.}
Which sentence from a student's analysis does the most analytical work?
\ch{``By opening with her own mistake, the author disarms readers who expect a lecture.''}{``The author uses imagery in paragraph 2.''}{``The author uses many rhetorical devices to persuade.''}{``The author uses ethos, pathos, and logos.''}
\end{q}

\begin{q}{4}{1.4 Appeals}{An argument from timing (``now, before the vote'') is kairos.}
A writer insists her argument matters ``this week, before the council votes, and not a day later.'' She is emphasizing
\chh{an analogy}{ethos}{logos}{kairos}
\end{q}

\section{Fitting choices to the situation}

The central idea of the course: a writer's choices are good or bad only \emph{relative to a situation}. The same speech would fail with a different audience. Questions in this category ask how a choice \textbf{accommodates} the audience, or how it reflects the writer's purpose or context.

\begin{compare}{Choice}{Why a writer might make it}{Typical question wording}
Delaying the thesis & The audience is resistant; reasoning must come first. & ``The writer's decision to withhold the main claim until paragraph 4\ldots'' \\
Using ``we'' and ``our'' & Creates a shared identity with the audience. & ``The shift from `you' to `we' serves to\ldots'' \\
Referring to shared history & Grounds the argument in values the audience already holds. & ``The reference to the town's founding\ldots'' \\
Plain, conversational register & Reaches a general audience; signals sincerity. & ``The informal diction is appropriate because\ldots'' \\
Technical detail & Satisfies an expert or skeptical audience demanding proof. & ``The inclusion of measurements\ldots'' \\
\end{compare}

\begin{q}{3}{1.5 Accommodating audience}{``We'' builds shared identity between writer and audience.}
Midway through a speech, a mayor shifts from ``you must conserve water'' to ``we have always shared what little we had.'' The shift most likely serves to
\ch{blame residents for the shortage}{signal a change of topic}{present conservation as a shared tradition rather than an order}{introduce technical data}
\end{q}

\begin{q}{3}{1.5 Accommodating audience}{Shared values first, then the request, for a resistant audience.}
Writing to factory owners who oppose a safety law, a reformer begins with their pride in their businesses. This choice is effective mainly because it
\ch{avoids the subject of the law entirely}{proves the law is unnecessary}{establishes common ground before making a request the audience may resist}{appeals to readers who already support the law}
\end{q}

\section{Writing the rhetorical analysis essay (Q2)}

\subsection{The prompt}
Q2 gives a nonfiction passage (often a speech or letter) with an introductory paragraph and asks you to ``write an essay that analyzes the rhetorical choices [the writer] makes to [achieve a purpose].'' The prompt usually names the purpose for you. Use it.

\subsection{A thesis that earns Row A}
A thesis must make a \emph{defensible claim about the writer's choices and purpose}. Strong pattern: \textbf{Writer + choices (described, not just named) + purpose + audience}.

\begin{study}{Thesis upgrade}
Weak: ``Smith uses ethos, pathos, and logos to persuade her audience.'' (names labels only; no analysis.)
Strong: ``Smith first wins over a skeptical board by honoring its caution, then uses the story of one student's commute to make the cost of inaction concrete, so that by the end, approving the bus route feels like the cautious choice.''
\end{study}

\subsection{Body paragraphs that earn Row B}
\begin{enumerate}[leftmargin=1.5em]
\item Organize by \textbf{choice} or by the \textbf{movement} of the text (opening, turn, conclusion), not by appeal.
\item For each choice: quote briefly, then explain \emph{how} it works on \emph{this} audience and \emph{why} it serves the purpose.
\item Use precise verbs: \emph{concedes, reframes, dignifies, unsettles, invites, undercuts, juxtaposes, amplifies}.
\item Avoid summary. If a sentence could appear in a book report, cut it.
\end{enumerate}

\subsection{Reaching Row C}
Explain the \textbf{significance} of the choices given the full situation (why this approach was needed for this audience at this moment), trace how the writer's strategy \textbf{develops or shifts}, or examine a \textbf{tension} the writer must manage (honoring an audience while criticizing it).

\begin{q}{1}{1.6 RA essay}{Labels alone do not make a defensible analytical thesis.}
Which thesis would most likely \emph{fail} to earn the Row A point on a rhetorical analysis essay?
\ch{``The author uses rhetorical devices such as ethos, pathos, and logos in his speech.''}{``King answers the clergy's charge of impatience by reframing waiting as its own form of harm.''}{``By narrating her own failure first, the speaker earns the right to advise graduates.''}{``The governor's plain diction and shared memories turn a disaster report into a call for unity.''}
\end{q}

\begin{q}{2}{1.6 RA essay}{Commentary explains how the choice affects this audience and serves the purpose.}
After quoting a writer's vivid description of a flooded classroom, a student should next
\ch{name the device as imagery and move on}{explain how the image makes the risk real for readers who have never seen a flood, advancing the writer's call for funding}{quote another vivid description}{summarize the rest of the paragraph}
\end{q}

\begin{q}{1}{1.6 RA essay}{Organize by choice or movement, not by appeal.}
Which organization is usually strongest for a rhetorical analysis essay?
\ch{By the writer's major choices, in the order they work on the audience}{One paragraph each for ethos, pathos, and logos}{A summary paragraph followed by a list of devices}{One paragraph per sentence of the passage}
\end{q}

\makekey
\end{document}
